\documentclass[10pt,a4paper]{amsart}
\usepackage[margin=19mm]{geometry}
\usepackage{amsmath,amssymb,mathtools}
\usepackage[scr=rsfs,cal=euler]{mathalfa}
\usepackage{xfrac}
\usepackage[unicode,hidelinks,bookmarks=false]{hyperref}
\newcommand{\ie}{i.e.\ }
\newcommand{\cf}{cf.\ }
\newcommand{\resp}{respectively\ }
\newcommand{\Subsection}{\subsection}
\providecommand{\nobreakdash}{\nobreak\hbox{-}}
\setlength{\emergencystretch}{2em}
\multlinegap0pt
\usepackage{amssymb}
\usepackage{dsfont}
\usepackage{enumerate}
\usepackage{csquotes}
\usepackage{mathtools}
\newtheorem{lemma}{Lemma}
\title{Waring–Goldbach problems for one square and higher powers}
\author{Geovane Matheus Lemes Andrade}
\date{Source: arXiv:2603.05550v1 (CC BY 4.0)}
\begin{document}
\maketitle

\noindent\emph{Source and licence.} Waring–Goldbach problems for one square and higher powers. Version arXiv:2603.05550v1 (CC BY 4.0), distributed by arXiv under the Creative Commons Attribution 4.0 International licence, http://creativecommons.org/licenses/by/4.0/, as recorded in the arXiv licence metadata for this exact version and shown on the abstract page https://arxiv.org/abs/2603.05550v1. Full licence notice: \texttt{licenses/arxiv-2603.05550-CC-BY-4.0.txt}.\par\smallskip

\noindent\textit{Editorial scope.} Counted unit: the article's lemma epsilon (source0.tex lines 309--312). The article's complete original proof of it is delivered with it (source0.tex lines 314--321, one \texttt{proof} environment of 8 lines). No mathematics is written, repaired or re-derived: the statement and the proof are the article's own lines, in the article's own notation, and no retained line is substituted or re-flowed.  Retained uncounted as the source-local context the proof needs: lines 293--295; lines 297--299.  Nothing was omitted from any retained span, so the package carries no omission record.  No retained line contains a citation, so the wrapper carries no bibliography and no bibliography entry.  No retained line contains a figure environment, an image-inclusion command or drawing code, so no image is delivered and the package declares no asset dependency.  Editorial formatting only, in addition: an AMS \texttt{amsart} preamble that restates the article's own declarations and macros used by the retained lines, namely the environments \texttt{lemma} on separate counters, together with the standard packages those lines need.  The retained environments print in this document as Lemma 1 in order of appearance, whereas the article prints the same items with the numbering of its own full document.  The complete native source of the article as distributed in the arXiv e-print for this exact version is bundled unchanged as \texttt{sources/195809/source0.tex}.
\begin{align}\label{j1j2}
    J_1 \ll n^{\frac{6}{5}+\epsilon}, \qquad J_2 \ll n^{1+\epsilon}.
\end{align}

\begin{align}\label{j3j4}
    J_3 &\ll n^{1-\lambda+\epsilon},
\end{align}

\begin{lemma}\label{epsilon}
   For $1\leq j \leq 2$, one has 
   $$ \int \limits_{\mathcal{E}_j} F_j(\alpha)e(-\alpha n)d\alpha \ll n^{\Theta_j - \tau}.$$
\end{lemma}

\begin{proof}
    By H\"older’s inequality, we deduce that 
    \begin{align*}
        \int \limits_{\mathcal{E}_1} |F_1(\alpha)|\,d\alpha &\leq J_1^{1/2}J_3^{1/2} \sup \limits_{\alpha \in \mathcal{E}_1}|f_{5}(\alpha)|,\\
        \int \limits_{\mathcal{E}_2} |F_2(\alpha)|\,d\alpha &\leq J_1^{1/4}J_2^{1/4}J_3^{1/2} \sup \limits_{\alpha \in \mathcal{E}_2}|f_{5}(\alpha)|^{1/2}.
    \end{align*}
    Employing the bounds in \eqref{j1j2} and \eqref{j3j4}, together with the definition of $\mathcal{E}_j$ and a straightforward calculation, the desired estimate follows.
\end{proof}

\end{document}
